\chapter{Aortic Aneurysm}
\index{Aortic Aneurysm}

\section*{Definition and Overview}
An aortic aneurysm is an abnormal, balloon-like bulge in the wall of the aorta, the largest artery in the body, which carries blood from the heart through the chest and abdomen. The wall of an aneurysm is weakened and thinned, so it can gradually enlarge over years and, if it grows too large, may tear or burst (rupture) -- a catastrophic, life-threatening emergency.

Most aneurysms occur in the abdominal portion of the aorta (abdominal aortic aneurysm, or AAA); others occur in the chest (thoracic aortic aneurysm). The two behave differently in terms of their causes, risk factors, and the size at which repair is recommended. Many aortic aneurysms cause no symptoms at all and are discovered incidentally during imaging for other conditions, which is why screening programmes and surveillance are so important.

\section*{Epidemiology}
Abdominal aortic aneurysms affect roughly one in twenty men over the age of 65 in screened populations and are several times more common in men than in women. The risk rises steeply with age, and smoking is the strongest modifiable risk factor; the combination of smoking, hypertension, and older age accounts for most cases. Thoracic aortic aneurysms are less common and occur at younger ages when associated with inherited conditions.

A family history of aneurysm, high blood pressure, and connective tissue disorders such as Marfan syndrome all increase risk. Because the condition is usually silent and rupture is often fatal, several countries operate ultrasound screening programmes for men at certain ages, and similar selective screening is offered to women with risk factors.

\section*{Aetiology and Pathophysiology}
The aortic wall normally remains strong through a balance of structural proteins -- particularly elastin and collagen -- and ongoing cellular repair. An aneurysm forms when this balance fails and the wall progressively weakens and dilates:

\begin{itemize}
    \item \textbf{Atherosclerosis:} fatty infiltration and chronic inflammation weaken the media of the wall; this is the dominant cause of abdominal aneurysms
    \item \textbf{High blood pressure:} sustained pressure stretches the wall and accelerates its degeneration
    \item \textbf{Smoking:} directly damages the arterial wall, increases inflammatory degradation, and speeds up aneurysm growth
    \item \textbf{Connective tissue disorders:} conditions such as Marfan syndrome, Ehlers--Danlos syndromes, and Loeys--Dietz syndrome cause abnormal structural proteins and typically produce thoracic and root aneurysms
    \item \textbf{Bicuspid aortic valve:} an inherited valve difference associated with enlargement of the ascending aorta
    \item \textbf{Family history:} about a fifth of people with AAA have an affected family member
    \item \textbf{Inflammation and infection:} rare causes including aortitis from vasculitis, or mycotic aneurysm from bacterial infection
    \item \textbf{Trauma:} blunt chest or abdominal injury can damage the aortic wall
\end{itemize}

\section*{Clinical Features}
\begin{itemize}
    \item \textbf{Usually none:} most aneurysms are silent until they are large or rupture
    \item \textbf{A pulsating feeling:} a palpable, throbbing mass in the abdomen, sometimes noticed by the person or a doctor
    \item \textbf{Deep, constant pain:} abdominal, back, or flank pain, especially with a large or expanding aneurysm
    \item \textbf{Thoracic symptoms:} chest or upper back pain, hoarseness, cough, or difficulty swallowing from compression of nearby structures
    \item \textbf{Rupture:} sudden, severe, tearing pain in the abdomen, chest, or back, with collapse, sweating, nausea, and a rapid, weak pulse
    \item \textbf{Blood loss:} dizziness, fainting, pallor, and reduced urine output as internal bleeding progresses
\end{itemize}

\section*{Differential Diagnosis}
\begin{itemize}
    \item \textbf{Other causes of abdominal pain:} pancreatitis, renal colic, peptic ulcer disease, diverticulitis, gallstones, or bowel obstruction
    \item \textbf{Musculoskeletal back pain:} disc disease, muscular strain, or spinal stenosis
    \item \textbf{Cardiac chest pain:} myocardial infarction or angina
    \item \textbf{Other vascular emergencies:} aortic dissection, mesenteric ischaemia, or a ruptured spleen
    \item \textbf{Abdominal or thoracic masses:} tumours, cysts, or enlarged lymph nodes
    \item \textbf{Aortic dissection:} which produces similar severe pain and may coexist with or be misdiagnosed as an aneurysm
\end{itemize}

\section*{When to Seek Medical Care}
See a doctor if you have risk factors -- smoking, high blood pressure, or a family history of aneurysm -- and develop unexplained abdominal or back pain. If you know you have an aneurysm, attend all surveillance appointments so that growth can be monitored and surgery timed appropriately.

\textbf{Contact your local emergency services for:}
\begin{itemize}
    \item Sudden, severe, or tearing pain in the abdomen, chest, or back
    \item Collapse, fainting, or dizziness occurring with the pain
    \item A rapid, weak pulse, sweating, or pale, clammy skin
    \item Any sudden, severe symptom that you are worried about -- rupture is a time-critical emergency
\end{itemize}

\section*{Diagnosis and Investigations}
\begin{itemize}
    \item \textbf{Physical examination:} a pulsatile abdominal mass may be felt in thinner people, but its absence does not rule out an aneurysm
    \item \textbf{Ultrasound:} the standard first-line test for the abdomen, and the basis of screening and surveillance programmes
    \item \textbf{CT angiography:} the most detailed test, used to measure the aneurysm precisely, assess anatomy, and plan open or endovascular repair
    \item \textbf{MRI or MR angiography:} an alternative for detailed imaging, particularly of the thoracic aorta, without ionising radiation
    \item \textbf{Blood pressure and cardiovascular risk assessment:} blood pressure measurement, cholesterol, and renal function to guide risk-factor management
    \item \textbf{Genetic assessment:} where a connective tissue disorder or strong family history is suspected
    \item \textbf{Chest X-ray or echocardiography:} may suggest a thoracic aneurysm, which is then defined with CT or MRI
\end{itemize}

\section*{Management}
\begin{itemize}
    \item \textbf{Surveillance:} small aneurysms are monitored with regular ultrasound or CT; the interval depends on size and growth rate
    \item \textbf{Blood pressure control:} antihypertensive therapy reduces wall stress and slows expansion
    \item \textbf{Smoking cessation:} stopping smoking slows growth and lowers rupture risk, and is the single most effective lifestyle measure
    \item \textbf{Healthy lifestyle:} balanced diet, lipid management, and regular, moderate exercise within the limits advised by the clinician
    \item \textbf{Elective repair:} open surgery or endovascular aneurysm repair (EVAR) is offered when the aneurysm reaches a size at which the risk of rupture outweighs the risks of the procedure (typically around 5.5 cm for abdominal aneurysms, with lower thresholds in women)
    \item \textbf{Emergency repair:} immediate surgery or endovascular stenting for a ruptured or rapidly leaking aneurysm, after resuscitation
\end{itemize}

\section*{Complications}
\begin{itemize}
    \item \textbf{Rupture:} a catastrophic internal bleed with a high death rate even with emergency surgery
    \item \textbf{Aortic dissection:} a tear in the wall that can coexist with or complicate an aneurysm
    \item \textbf{Thromboembolism:} clot formation within the aneurysm that can break off and block arteries to the legs, kidneys, or bowel
    \item \textbf{Compression:} a large thoracic aneurysm pressing on the oesophagus, airways, or recurrent laryngeal nerve
    \item \textbf{Complications of repair:} bleeding, infection, kidney injury, and leakage or endoleak around an endovascular stent graft
    \item \textbf{Cardiac complications:} ischaemic heart disease risk is high in this population, both before and after repair
\end{itemize}

\section*{Prognosis and Outcomes}
The outlook depends mainly on the size and growth rate of the aneurysm. Small aneurysms can remain stable for years under surveillance, and many never reach a size requiring repair. The risk of rupture rises steeply as the aneurysm enlarges, which is why surgery is timed carefully rather than performed at diagnosis.

Elective repair is generally safe and effective, with a low operative mortality in well-selected patients and durable long-term results. A ruptured aneurysm is a different matter: many people die before reaching hospital, and even with emergency surgery the death rate remains high. Early detection through screening, along with blood pressure control and smoking cessation, saves lives.

\section*{Prevention and Self-Care}
\begin{itemize}
    \item Never smoke, and seek help to stop if you do -- it is the single most important step
    \item Control high blood pressure and cholesterol under a doctor's guidance
    \item Eat a balanced diet, keep physically active, and maintain a healthy weight
    \item Attend ultrasound screening if you are in an age and risk group offered it
    \item If you have a connective tissue disorder or a family history of aneurysm, discuss monitoring with your doctor
    \item Know the warning signs of rupture and contact local emergency services immediately if they occur
    \item Attend all scheduled surveillance imaging so that any growth is detected promptly
\end{itemize}

\section*{Sources and Further Reading}
The following organisations provide reliable, up-to-date information on aortic aneurysms for patients, students, and health professionals.

\begin{itemize}
    \item NHS. \textit{Information on abdominal aortic aneurysm, screening, and treatment.} https://www.nhs.uk/conditions/
    \item Mayo Clinic. \textit{Overview of aortic aneurysm, symptoms, and treatment.} https://www.mayoclinic.org/
    \item NICE. \textit{Guidance on abdominal aortic aneurysm diagnosis and management.} https://www.nice.org.uk/
    \item MedlinePlus. \textit{Health information on aortic aneurysm and screening.} https://medlineplus.gov/
    \item MSD Manual. \textit{Professional overview of aortic aneurysm and dissection.} https://www.msdmanuals.com/
    \item World Health Organization. \textit{Resources on cardiovascular disease prevention.} https://www.who.int/
\end{itemize}
